\section{Methodology}\label{sec:method}

\subsection{The unlearning problem}

Machine unlearning arises when part of a model's training data must
be deleted after training. Let $D$ be the training data and $A$ the
training procedure, including the architecture, the hyperparameters,
and the random seed. The \emph{forget set} $F \subset D$ holds the
observations to be deleted, and the \emph{retain set} $D \setminus F$
holds the rest. The model trained on everything, $A(D)$, is the
\emph{reference}. The model retrained from scratch on the retain
set, $A(D \setminus F)$, is the \emph{oracle}, or
\emph{counterfactual}, the ground truth of deletion. Because
retraining is often unaffordable, \emph{approximate unlearning}
instead edits the reference in place with an \emph{unlearning
operator} $U$, aiming for a model that a prescribed audit cannot
distinguish from the oracle.

\begin{table}[t]
\centering
\caption{Deletion levels for a requested period $T$. Each level
includes all restrictions of the preceding level. \Lone{} is the
paper-wide default.}
\label{tab:deletion_levels}
\footnotesize
\begin{tabular}{@{}p{0.13\columnwidth}p{0.78\columnwidth}@{}}
\toprule
Level & Prohibited use of $T$ and representative examples \\
\midrule
\Lzero{} &
\textbf{Target deletion.}
Use in forecast targets, such as prediction labels. \\
\Lone{} &
\textbf{Window deletion.}
Direct appearance in inputs or targets, such as context windows,
and use in the normalization statistics that scale them. \\
\Ltwo{} &
\textbf{Derived-data deletion.}
Dependence through feature construction, such as rolling statistics
and imputation. \\
\Lthree{} &
\textbf{Full-pipeline deletion.}
Influence on model-development decisions, such as hyperparameter
tuning and model selection. \\
\bottomrule
\end{tabular}
\end{table}

\subsection{A membership lattice for temporal deletion}

In time series, a deletion request does not name a forget set. It
names a period $T$, a set of dates whose use is no longer permitted.
A single date inside $T$ can touch a rolling forecaster in four
distinct ways, and each way is one level of
Table~\ref{tab:deletion_levels}. Consider one trading day in March
2020. It is the prediction \emph{target} of the windows that
forecast it (\Lzero). It sits inside the 64-day \emph{input context}
of windows whose targets fall up to three months later (\Lone). It
enters \emph{derived quantities} computed from it, such as rolling
features with long tails (\Ltwo). And it may have influenced
\emph{model-development decisions} such as hyperparameter tuning
(\Lthree). Each level forbids everything the preceding level does
plus one further form of dependence, so the four readings are
nested. We call this ordering the membership lattice. The first
three levels change only which training examples are retained,
giving nested forget sets $F_\ell(T)$, whose exact construction is
specified in the released protocol.

Strictness has a measurable price in two currencies. The first is
data. Because windows overlap, \Lone{} removes windows whose target
dates lie outside $T$, and on this panel it surrenders about $25\%$
more windows than \Lzero. The second is the measurement itself. The
level changes the forget set, and with it the measured memorization
gap (Section~\ref{sec:levels}). A temporal deletion request is therefore incomplete unless its
level is named. We adopt \Lone{} as the paper-wide default and use
\Lzero{} and \Ltwo{} as robustness checks. \Lthree{} may alter the
training procedure itself rather than only the retained examples.
It coincides with \Ltwo{} here by construction, because tuning and
model selection predate every deletable era.

At every level, the target of unlearning is the oracle retrained
without that level's forget set,
\begin{equation}
U\!\left(A(D),F_\ell(T)\right)
\;\approx_d\;
A\!\left(D\setminus F_\ell(T)\right),
\qquad
\ell\in\{0,1,2\},
\label{eq:unlearning}
\end{equation}
where $\approx_d$ is approximate equality in distribution under a
prescribed audit, and the method and the oracle use the same level.

% [PAGE-BUDGET L1] restore by removing the leading %% on this block
%\begin{figure}[t]
%\centering
%\begin{tikzpicture}[x=0.68cm, y=0.5cm, font=\footnotesize]
%  \fill[gray!20] (3.2,-3.95) rectangle (5.6,0.55);
%  \draw[gray!60] (3.2,-3.95) rectangle (5.6,0.55);
%  \node at (4.4,-0.05) {deleted era $T$};
%  \draw[decorate, decoration={brace, amplitude=3pt}]
%    (5.65,0.55) -- (6.95,0.55) node[midway, above=3pt] {context halo};
%  \draw[->] (0,-3.95) -- (8.3,-3.95);
%  \node[below] at (4.15,-3.95) {trading days};
%  \node[right] at (8.0,-0.15) {dropped at:};
%  \draw[fill=gray!8] (0.2,-1.10) rectangle (1.60,-0.65);
%  \fill[black] (1.60,-1.10) rectangle (1.78,-0.65);
%  \node[right] at (8.0,-0.875) {---};
%  \draw[fill=gray!8] (1.75,-1.75) rectangle (3.15,-1.30);
%  \fill[black] (3.15,-1.75) rectangle (3.33,-1.30);
%  \node[right] at (8.0,-1.525) {\Lzero{}, \Lone};
%  \draw[fill=gray!8] (3.50,-2.40) rectangle (4.90,-1.95);
%  \fill[black] (4.90,-2.40) rectangle (5.08,-1.95);
%  \node[right] at (8.0,-2.175) {\Lzero{}, \Lone};
%  \draw[fill=gray!8, very thick] (4.70,-3.05) rectangle (6.10,-2.60);
%  \fill[black] (6.10,-3.05) rectangle (6.28,-2.60);
%  \node[right] at (8.0,-2.825) {\Lone{} only};
%  \draw[fill=gray!8] (6.30,-3.70) rectangle (7.70,-3.25);
%  \fill[black] (7.70,-3.70) rectangle (7.88,-3.25);
%  \node[right] at (8.0,-3.525) {---};
%\end{tikzpicture}
%\caption{One deletion request, four readings. Each row is a training
%window: bar $=$ 64-day context, dark tick $=$ 5-day forecast target.
%\Lzero{} drops windows whose \emph{target} touches $T$ (rows 2--3).
%\Lone{} also drops windows whose \emph{context} touches $T$ (row 4,
%bold). Row 5 lies beyond the halo and survives both.}
%\label{fig:lattice}
%\end{figure}

\subsection{Counterfactuals and governance metrics}

For each deletion request and level, the oracle retrains the model from
scratch on the retain set under the reference model's training
protocol. An approximate method is judged by how closely it reproduces
this counterfactual, not by how much it degrades predictions on the
deleted period.

Let $\text{ref} = A(D)$ denote the reference model,
$\text{ora} = A(D \setminus F_\ell(T))$ the retrained counterfactual,
and $m = U(\text{ref}, F_\ell(T))$ the model produced by an approximate
operator. The forget-set loss $\mathcal{L}_F(\cdot)$ measures how
well a model still fits the deleted data, so a higher loss indicates
more forgetting. The test score $s_{\text{test}}(\cdot)$
measures forecast quality on the test set, oriented so that higher is
better.

\textbf{Memorization gap $G_F$.} The memorization gap is the reference
model's excess fit over the counterfactual on the forget set,
\begin{equation}
G_F = \mathcal{L}_F(\text{ora}) - \mathcal{L}_F(\text{ref}).
\label{eq:gap}
\end{equation}
A gap near zero means the forget set was barely memorized; a large gap
means it was memorized strongly. $G_F$ is the era-level analogue of the per-example memorization
score of Feldman and Zhang~\cite{feldman2020what}. It carries the
units of the target a model was fit to, and is rescaled as
$G_F/\mathrm{Var}(y_T)$ when a model with a different target is
compared.

\textbf{Closure $\gapclosed$.} The closure is the fraction of the
memorization gap that a method $m$ removes,
\begin{equation}
\gapclosed(m) =
\frac{\mathcal{L}_F(m) - \mathcal{L}_F(\text{ref})}{G_F}.
\label{eq:kappa}
\end{equation}
A closure of $0$ means the method forgets nothing and remains at the
reference, $1$ means it reaches the counterfactual, and above $1$
means it over-forgets.

\textbf{Deletion effect $\Delta s_{\text{test}}$.} The deletion effect
is the change in test score caused by exact deletion, measured one day
at a time so that offsetting daily differences are not hidden by the
average. The reference and the counterfactual score the same
cross-section on each evaluation day $t$, so the daily change is
$\Delta s_t = s_t(\text{ora}) - s_t(\text{ref})$, and
\begin{equation}
\Delta s_{\text{test}} =
\frac{1}{N_{\text{test}}}\sum_{t=1}^{N_{\text{test}}} \Delta s_t .
\label{eq:delta}
\end{equation}
A positive $\Delta s_{\text{test}}$ means deletion improves the
test score, and a negative value means the deleted period carried
useful signal. Where no confusion arises we write $\Delta$ for
$\Delta s_{\text{test}}$ and $G$ for $G_F$.

\subsection{Auditing deletion}

The goal of an audit is to determine whether information from the
deleted period remains detectable in the model. Throughout,
\emph{verified} means one thing only: detectors of stated power found
no gap above their minimum detectable effect. It never means proof
that no trace remains.

The audit forms two populations by retraining under different random
seeds: references that include $T$ in training, and counterfactuals
that omit it. This is the access model a model owner or an appointed
auditor can realistically arrange: fixed code, fixed data cuts, and
multiple retraining seeds. Given a per-model statistic $a(\cdot)$,
the audit separation is
\begin{equation}
d' = \frac{\bigl\lvert\, \overline{a(\text{ora})} -
\overline{a(\text{ref})} \,\bigr\rvert}{\mathrm{sd}_{\text{pool}}},
\label{eq:dprime}
\end{equation}
the standardized distance between the two population means, where
$\mathrm{sd}_{\text{pool}}$ is the pooled standard deviation across
seeds. A large $d'$ means the populations separate and deletion is
detectable; $d' \approx 0$ means they are indistinguishable and
deletion cannot be certified from $a$. We take $a = \mathcal{L}_F$
because it is evaluated on the deleted data, where the populations
separate sharply. We report $d'$ rather than an AUC because $d'$
does not saturate.

The granularity of $a$ matters. Aggregating over the whole deleted
period gives an era-level audit, and scoring individual overlapping
windows gives a window-level audit.
